\abstract{
\textbf{Covariance-Aware Policy Optimization (CAPO)} is a GRPO-style, variance-aware trajectory reweighting framework for Reinforcement Learning with Verifiable Rewards (RLVR). CAPO starts from a generalized Bayesian working model for per-trajectory policy-gradient statistics and replaces token-independence assumptions with a principled treatment of within-trajectory covariance. With a Jeffreys prior on the global scale, the posterior-precision rule recovers $\Delta L$ normalization as the correctly specified independence special case. Beyond independence, CAPO’s theory shows that structured token dependence induces an \emph{induced variance law} $v(L)=\mathrm{Var}(g\mid L)$ that may grow linearly (summable/short-range dependence) or superlinearly (persistent/long-range dependence), and that the statistically efficient GRPO group reweighting is inverse-variance.

To connect theory to practice, we express CAPO as a \emph{modular plug-in algorithm}: the policy update requires only an estimator of $v(L)$ (up to scale), which is converted into within-group precision weights. We develop and analyze four plug-ins of increasing structure: \textbf{L-CAPO}, the length-only/independence plug-in (including $\Delta L$ as the fixed-$\beta{=}1$ case); \textbf{LV-CAPO}, an empirical-Bayes plug-in that jointly estimates a length exponent and a structured dependence shape; \textbf{CAPO-Q}, a quadratic variance-law plug-in motivated by compound-symmetry/random-intercept covariance; and \textbf{CAPO-HAC}, a long-run-variance plug-in based on lag-window (Newey--West) estimation. We provide closed-form induced-variance derivations, grouped (prompt-wise) EB objectives and gradients, robustness guarantees for plug-in inverse-variance weighting, and diagnostics for when CAPO reduces variance without introducing bias. Across synthetic covariance benchmarks and CountDown RLVR experiments, CAPO variants stabilize training and improve efficiency relative to $\Delta L$ and GRPO/DAPO/Dr.~GRPO normalizations.

\textbf{Keywords}: Reinforcement Learning with Verifiable Rewards (RLVR); Bayesian Methods; Variance Reduction; Policy Optimization; Large Language Models; Sequence Modeling; Stochastic Optimization; Robust Training; Statistical Inference.

}
